% =========================
% Explanation: Statistical Analysis Plan (informal, readable)
% =========================
\subsection*{Explanation of Statistical Analysis Plan (Pre-specified) (informal)}
This section explains how we analyze results across stochastic runs and how we avoid common pitfalls:
overstating significance, ignoring multiple comparisons, or drawing conclusions from too few seeds.
Because each optimizer run is random, we treat performance as a distribution and compare methods using robust summaries
and appropriate statistical tests.

\begin{itemize}
    \item \textbf{Why pre-specify the plan?}
    Benchmarking involves many degrees of freedom (which budgets to show, which metrics to emphasize, which baselines to compare,
    and which seeds to include). Pre-specifying the analysis plan reduces ``researcher degrees of freedom'' and increases credibility,
    because it prevents selecting tests or metrics after seeing the outcomes.

    \item \textbf{Aggregation across repetitions.}
    For each (method, benchmark instance, budget) we run \(R\) independent repetitions (different optimizer seeds).
    We report robust summaries:
    \begin{itemize}
        \item \emph{Primary:} median with interquartile range (IQR), because distributions can be skewed and may contain occasional failures.
        \item \emph{Secondary:} mean \(\pm\) standard deviation, to support comparisons with prior work.
    \end{itemize}
    For traces (e.g., best-so-far vs.\ evaluations), we aggregate pointwise medians with IQR bands.

    \item \textbf{Handling right-censoring for fixed-target metrics.}
    Fixed-target metrics (e.g., time to reach a threshold) can be \(\infty\) if a run never reaches the target within budget.
    Such outcomes are \emph{right-censored}.
    We handle this explicitly by:
    (i) reporting Kaplan--Meier survival curves (probability of not having reached the target by evaluation \(t\)),
    and/or (ii) using a conservative imputation such as \(T = B_{\max}+1\) for rank-based tests, with this choice stated clearly.

    \item \textbf{Inferential tests for stochastic optimizers.}
    Because performance distributions are rarely Gaussian and sample sizes per benchmark are modest, we prefer nonparametric tests.
    Typical choices include:
    \begin{itemize}
        \item \emph{Paired Wilcoxon signed-rank test} when runs can be paired by shared seeds or common random numbers
        (reduces variance and increases power).
        \item \emph{Mann--Whitney U test} when pairing is not justified (independent samples).
        \item \emph{Permutation tests} as a robust alternative when assumptions are unclear.
    \end{itemize}
    Tests are applied to per-run scalar outcomes (e.g., regret or HV at budget \(B\), or time-to-target).

    \item \textbf{Multiple comparisons control.}
    Benchmarking typically compares many methods across multiple datasets, budgets, and metrics.
    Without correction, false positives are likely.
    We therefore control family-wise error rate using Holm--Bonferroni (or control FDR via Benjamini--Hochberg if explicitly chosen),
    and report corrected \(p\)-values or adjusted significance thresholds.

    \item \textbf{Effect sizes (practical significance).}
    Statistical significance alone does not imply a meaningful improvement.
    We therefore report effect sizes alongside \(p\)-values, such as:
    \begin{itemize}
        \item Vargha--Delaney \(\hat{A}_{12}\) (probability that method A outperforms method B on a randomly chosen run),
        \item Cliff's delta (magnitude and direction of stochastic dominance),
        \item or standardized median differences where appropriate.
    \end{itemize}
    These quantify practical impact and robustness.

    \item \textbf{Outliers, failed runs, and invalid configurations.}
    HPO/NAS frequently produces failed evaluations (OOM, NaNs, divergence, invalid architectures).
    To avoid hidden bias, we specify a uniform handling rule:
    failed/invalid trials are assigned a penalized objective value (worse than any valid value on that benchmark)
    and still count toward the evaluation budget.
    We report failure frequencies and causes. We do not discard outliers unless there is a documented logging/evaluation error.

    \item \textbf{Across-benchmark summarization.}
    When summarizing across benchmark instances, we avoid pooling raw objective values across incomparable scales.
    Instead, we use (i) per-instance normalization (e.g., regret normalized by an oracle gap) and then (ii) macro-averaging across instances.
    We also report performance profiles / ECDFs to show how often a method is within a given factor of the best.

    \item \textbf{Ablations and sensitivity analysis.}
    To attribute improvements to specific components, we pre-specify ablations:
    \begin{itemize}
        \item remove the ``Net'' adaptations (vanilla backbone only),
        \item remove each novel component one at a time,
        \item disable diversity control (tests mode-collapse sensitivity),
        \item swap constraint handling mechanism (repair vs.\ rejection) while keeping accounting fixed,
        \item vary key method hyperparameters (e.g., elite fraction, smoothing, mixture rate).
    \end{itemize}
    Ablations are run under the same budgets, seeds, and evaluation pipeline to ensure clean attribution.

    \item \textbf{What constitutes evidence.}
    We consider the method improved if it shows:
    (i) consistent improvements in the primary metric across benchmark instances and seeds,
    (ii) nontrivial effect sizes (not only tiny significant differences),
    and (iii) robustness to reasonable hyperparameter perturbations (sensitivity analysis).
\end{itemize}
